\subsection{Overview}
CPUID is the architectural discovery instruction. It exposes base-profile identity, optional
architectural extensions, and program-visible tuning properties. It is not a dump of internal
microarchitectural structures.

A leaf describes one meaningful information group. The result index selects one Q-sized result
from that group.

\subsection{Query Model}
The instruction form is @SYNTAX@. The selected D register contains a 64-bit query selector
before execution and receives the selected 64-bit result after execution. The selector is divided into
a 32-bit class, a 16-bit leaf within that class, and a 16-bit result index.

@CALLING_CONVENTION_TABLE@

@SELECTOR_TABLE@

\subsection{Discovery Classes}
\noindent\textbf{Base profile.} The base Bedrock ISA profile contains ordinary integer, data movement,
control-flow, atomic, data-register banking, cache/TLB, segmentation, paging, control-register, and
repeat-prefix instruction groups. Instruction groups with their own optional-extension leaves, such as
virtualization acceleration, are not implied by the base profile.

\noindent\textbf{Floating-point.} Floating-point discovery is intentionally coarse. An implementation reports
whether the base floating-point extension is available, and separately whether the transcendental
floating-point group is available. Double precision, fused multiply-add/subtract, floating-point
conditional move, conversion, classification, and ordinary floating-point memory/register forms are base
features of the floating-point extension.

\noindent\textbf{Data-register banking.} Data-register banking instructions are part of the base instruction set.
DBANK implementation-property discovery reports the implemented bank count and selector width; object
attributes record the minimum bank count required by a linked image. Source-language ABIs still normalize
DBANK to zero at public boundaries.

\noindent\textbf{Virtualization acceleration.} ENCINST availability is discovered through the optional-extension
class. The instruction reference groups ENCINST under Virtualization Acceleration because it supports
canonical binary emission for JITs, hypervisors, and binary translators rather than ordinary application
execution.

\noindent\textbf{Unavailable optional groups.} If an optional instruction group is not reported by CPUID,
executing an instruction from that group raises \texttt{ILLEGAL\_INSTRUCTION}. Extension availability uses the
ordinary instruction-encoding fault policy.

\noindent\textbf{Implementation properties.} Implementation-property leaves expose only information a program can
reasonably use to select code paths or tune algorithms. Structures such as register renaming, reorder buffers,
issue queues, and physical register files are not architectural discovery bits. The initial execution-property
leaf exposes only \texttt{OUT\_OF\_ORDER}. Processor-topology discovery reports the total exposed hardware-thread
count and the current hardware-thread identity; it does not report the set of CPUs currently online under an
operating system.

@POLICY_TABLE@

@CLASS_TABLE@

\subsection{Leaf Directory}
@LEAF_TABLE@

@LEAF_DETAILS@

\subsection{Bit Fields}
@BIT_FIELD_BLOCKS@
